Una pequeña empresa artesanal fabrica dos tipos de jabones ecológicos: de \textbf{lavanda}  y de \textbf{romero}. Ambos se elaboran a partir de aceites esenciales y glicerina vegetal, y el objetivo de la empresa es \textbf{maximizar el beneficio semanal}.

Cada kilogramo de jabón requiere determinadas cantidades de \textbf{aceite esencial (ml)} y \textbf{glicerina (ml)}, y genera un beneficio unitario según la siguiente tabla:

\[
\begin{array}{lccc}
\text{Producto} & \text{Aceite esencial (ml/kg)} & \text{Glicerina (ml/kg)} & \text{Beneficio (€ / kg)} \\
\hline
\text{Jabón de lavanda } & 40 & 100 & 5 \\
\text{Jabón de romero }  & 60 & 80  & 4 \\
\hline
\end{array}
\]

La empresa dispone semanalmente de \textbf{8\,000 ml de aceite esencial} y \textbf{12\,000 ml de glicerina}. 

Además, por cuestiones de marketing, se deben fabricar al menos \textbf{60 kg en total} de jabón.

Por último, como máximo el 50\% de la glicerina consumida debe destinarse al jabón de romero.

El siguiente modelo de programación permite obtener el plan de producción óptimo, donde ($x_1$ y $x_2$ representan los kg de jabón de lavanda y de romero, respectivamente):

\[
\begin{aligned}
\mbox{max.} \quad & z = 5x_1 + 4x_2 \\
\text{s.a.} \quad 
& 40x_1 + 60x_2 \le 8\,000 & & \text{(disponibilidad de aceite esencial)} \\
& 100x_1 + 80x_2 \le 12\,000 & & \text{(disponibilidad de glicerina)} \\
& x_1 + x_2 \ge 60 & & \text{(demanda mínima total)} \\
& 5x_1 - 4x_2 \ge 0 & & \text{(glicerina del romero)} \\
& x_1, x_2 \ge 0
\end{aligned}
\]

Se conoce la tabla correspondiente a la solución óptima del problema:

% \begin{center}
% \begin{tabular}{c|c|cccccc|}
%  & $z$ & $x_1$ & $x_2$ & $h_1$ & $h_2$ & $h_3$ & $h_4$\\ 
%        & $-600$ & $0$ & $0$ & $0$ & $-1/20$ & $0$ & $0$ \\ 
% \hline
% $h_1$ & 1520  & $0$ & $0$ & $1$ & $-27/50$ & $0$ & $-14/5$\\ 
% $h_3$ & 72  & $0$ & $0$ & $0$ & $11/1000$ & $1$ & $1/50$\\ 
% $x_2$ & 60  & $0$ & $1$ & $0$ & $1/200$ & $0$ & $1/10$\\ 
% $x_1$ & 72  & $1$ & $0$ & $0$ & $3/500$ & $0$ & $-2/25$\\ 
% \hline
% \end{tabular}
% \end{center}


\begin{center}
\begin{tabular}{c|c|cccccc|}
 & $z$ & $x_1$ & $x_2$ & $h_1$ & $h_2$ & $h_3$ & $h_4$\\ 
Phase 2 & $-600$ & $0$ & $0$ & $0$ & $-1/20$ & $0$ & $0$ \\ 
\hline
$h_1$ & 1100  & $0$ & $0$ & $1$ & $-23/40$ & $0$ & $-7/2$\\ 
$h_3$ & 75  & $0$ & $0$ & $0$ & $9/800$ & $1$ & $1/40$\\ 
$x_2$ & 75  & $0$ & $1$ & $0$ & $1/160$ & $0$ & $1/8$\\ 
$x_1$ & 60  & $1$ & $0$ & $0$ & $1/200$ & $0$ & $-1/10$\\ 
\hline
\end{tabular}
\end{center}




\begin{enumerate}[label=\alph*)]
    \item Explicar por qué el requisito ``como máximo el 50\% de la glicerina consumida debe destinarse al jabón de romero'' se formula como $5x_1 - 4x_2 = 0$.

    \item Interpretar la solución óptima en términos del beneficio, el plan de producción y el cumplimiento de los requisitos.
    
    \item Identificar la base y la inversa de la base de la solución óptima.

    \item ¿Cuál sería el precio que debería estar dispuesta a aceptar la empresa para renunciar a 1 litro de glicerina en una semana dada?¿Cuántos litros estaría dispuesta a pagar por ese precio?
\end{enumerate}